\begin{theorem}[g-BQO]
For every nonidentity $g\in\mathsf{IncSeq}$ and every binary relation $R$
on a discrete type $A$,
\[
(A,R)\text{ is a $g$-better-relation}
\quad\Longleftrightarrow\quad
(A,R)\text{ is a better-relation}.
\]
In particular, if $Q$ is a quasi-order with relation $\leq$, then
\[
Q\text{ is $g$-bqo}\quad\Longleftrightarrow\quad Q\text{ is bqo}.
\]
\end{theorem}
\begin{proof}
Fix $g\neq\mathsf{IncSeqId}$.  By \noderef{MainProp}, choose continuous
homomorphisms in both directions between
$(\mathsf{IncSeq},\mathsf{RightComp}(g))$ and
$(\mathsf{IncSeq},\mathsf{ShiftMap})$.  We first record the transport
argument for the second clauses in \noderef{GBetterRel} and
\noderef{BetterRel}.  If $K$ is a continuous homomorphism from a source
system $(\mathsf{IncSeq},s)$ to $(\mathsf{IncSeq},t)$ and $H$ is a
continuous homomorphism from $(\mathsf{IncSeq},t)$ to
$(A,R^{\complement})$, then $H\circ K$ is locally constant and is a
homomorphism from $(\mathsf{IncSeq},s)$ to $(A,R^{\complement})$:
local constancy follows by first choosing a prefix deciding the required
finite output of $K$ and then a prefix deciding $H$ on that output, while
the homomorphism equation follows from
$K(s(x))=t(K(x))$ and $H(t(y))$ being related by $R^{\complement}$ to
$H(y)$.  The reverse source homomorphism gives the converse implication.
Thus the absence of the complement homomorphism for $g$ is equivalent to
the absence of the complement homomorphism for the ordinary shift.  The
source actions here are the maps defined in \noderef{RightComp} and
\noderef{ShiftMap}.

For completeness, the two clauses in \noderef{GBetterRel} are equivalent by
the standard finite-prefix fusion argument.  If the first clause fails for
a locally constant $\varphi$, recursively choose, after every finite
prefix $u$, an extension $v$ witnessing failure of the restricted
homomorphism.  At stage $n$ only finitely many values of $\varphi$ have
been queried, so local constancy lets us extend the current prefix beyond
all of them.  The increasing union of these prefixes defines an increasing
enumeration $x$ for every input and, by the choice at each stage, the map
$y\mapsto\varphi(x\circ y)$ is locally constant and satisfies
$\neg R(\varphi(x\circ y),\varphi(x\circ y\circ g))$ for every $y$.
This is a complement homomorphism, contradicting the second clause.
Conversely, suppose $H$ is a complement homomorphism.  For every $f$, the
restriction $x\mapsto H(f\circ x)$ is still locally constant and still
satisfies the complement relation, because the homomorphism equation for
$H$ can be evaluated at $f\circ x$.  It therefore cannot also be an
$R$-homomorphism: at any $x$ that would assert both $R(a,b)$ and
$\neg R(a,b)$.  Thus no restriction supplied by the first clause can
exist.  The same fusion uses $\mathsf{ShiftMap}$ in place of
$\mathsf{RightComp}(g)$ for \noderef{BetterRel}.  Hence the transported
second clauses also transport the first clauses, proving the first
displayed equivalence.

Suppose now that $r$ is a quasi-order on $Q$.  By \noderef{ContinuousRelHom}, a
homomorphism from the shift into the complement relation is precisely a
locally constant map $h:\mathsf{IncSeq}\to Q$ such that
\[
\neg r(h(x),h(\mathsf{ShiftMap}(x)))
\]
for every $x$.  By \noderef{BadMultiSequence} this says exactly that $h$
is a bad multi-sequence, and \noderef{Bqo} says that no such locally
constant map exists.  Therefore $\operatorname{BetterRel}(r)$ is
equivalent to $\operatorname{Bqo}(r)$.  Combining this with the first
equivalence gives the quasi-order specialization.
\end{proof}
